\section{Corpus}
\label{sec:corpus}

The corpus is the early learning standards documents of six US states: four
used for development (Arizona, California, Colorado, Texas) and two held out
entirely (Nevada, Kentucky). Table~\ref{tab:dataset-stats} summarizes what the
annotated portion of each contains; Appendix~\ref{sec:appendix-corpus} names
every source document and lists the pages retained at each tier.

Each document exists at up to three tiers, and the two reductions are not the
same kind of reduction. The \emph{full} tier is the published PDF as issued.
The \emph{trimmed} tier removes non-standards matter (front matter,
introductions, guiding-principles essays, appendices, acknowledgements) and
\textbf{retains every standard}; it reduces page count, and therefore cost, but
not coverage, and it is the tier the scale measurements of
\S\ref{sec:experiments-scale} and \S\ref{sec:experiments-scale-quality} run on.
Kentucky's 52 trimmed pages are all of Kentucky's standards. The \emph{subset}
tier is a genuine reduction in coverage: one to two domains per state, 8--15
pages each and 70 in total, chosen to be cheap to iterate on and to annotate
exhaustively. \textbf{Every headline quality number in this paper is computed
at this tier}; the one quality measurement at the trimmed tier
(\S\ref{sec:experiments-scale-quality}) is kept in its own table. Both cuts
were made by hand, which is favorable preprocessing; Table~\ref{tab:corpus-pages}
lists the retained pages so that either tier can be reconstructed from the
issuing agency's copy. It gives two numberings, because they are not the same
number: a page's index in the PDF file, and the page number printed on the
page, which differ by each document's unnumbered front matter. Kentucky's PDF
page 52 is printed page 39, and its subset also contains PDF page 65, whose
printed number is 52, so ``page 52'' is ambiguous on that document alone.
Printed numbers are recovered automatically and are reported only for the two
documents where that recovery is reliable across the whole tier.

Two sources need naming precisely. California's document is the state's
Preschool/Transitional Kindergarten Learning Foundations ``at a glance''
publication (68 pages), which was cut directly to its 13-page subset with no
trimmed tier. Colorado's is the 41-page Ages~3--5 volume of the state's
birth-to-eight guidelines; its trimmed tier is the whole document, so nothing
was cut from it, and its out-of-scope content is other age bands rather than
standards removed by trimming. The held-out canary is Nevada's 2023
pre-kindergarten standards, a multi-domain document; a separate Nevada
publication from 2025 covering only the social-emotional domain was collected
but is too narrow to support a generalization claim and is excluded from every
measurement (Table~\ref{tab:corpus-appendix}).

Two properties of Table~\ref{tab:dataset-stats} are results rather than
document facts. Its age-band column counts the distinct bands a state uses,
not coverage: every one of the 262 standards carries an age band. And
identifier collisions are zero in every state, which matters because the
identifier is the key under which a standard is stored (\S\ref{sec:schema}).

% AUTO-GENERATED by paper/analysis/generate_tables.py -- do not hand-edit.
% Source: paper/results/task8_20260904/dataset_stats.json
% Regenerate with: python paper/analysis/generate_tables.py
% Corpus tier (guardrail 1): every row is the _only_subset tier, never full-document.
\begin{table*}
  \centering
  \begin{tabular}{llrrrrrrrr}
    \hline
    \textbf{State} & \textbf{Role} & \textbf{Subset pp} & \textbf{Full pp} & \textbf{Dom.} & \textbf{Str.} & \textbf{Sub.} & \textbf{Standards} & \textbf{Age bands} & \textbf{ID coll.} \\
    \hline
    AZ & golden & 15 & 217 & 2 & 4 & 9 & 45 & 1 & 0 \\
    CA & golden & 13 & 68 & 3 & 7 & 18 & 94 & 3 & 0 \\
    CO & golden & 10 & 41$^{\dagger}$ & 3 & 10 & 0 & 48 & 1 & 0 \\
    TX & golden & 9 & 87 & 2 & 4 & 2 & 25 & 2 & 0 \\
    NV & held-out & 15 & 98 & 3 & 7 & 7 & 24 & 1 & 0 \\
    KY & held-out & 8 & 120 & 3 & 5 & 10 & 26 & 1 & 0 \\
    \hline
    \textbf{Total} & & 70 & & 16 & 37 & 46 & 262 & 6 & \textbf{0} \\
    \hline
  \end{tabular}
  \caption{Corpus descriptive statistics. Domain, strand and sub-strand columns count \emph{distinct} nodes in at least one standard's ancestor chain, so they are smaller than the detector's element counts in Table~\ref{tab:detector-headline}. \textbf{Age bands} is the state's distinct bands; its total is the union across states, not a column sum. CO's subsets derive from the 41pp ages-3--5 document ($\dagger$). All quality numbers in this paper are measured on the \emph{\_only\_subset} tier shown here (\S\ref{sec:corpus}). Source: \texttt{task8\_20260904}.}
  \label{tab:dataset-stats}
\end{table*}

